\documentclass[CJK,13pt]{beamer}
\usepackage{CJKutf8}
\usepackage{beamerthemesplit}
\usetheme{Malmoe}
\useoutertheme[footline=authortitle]{miniframes}
\usepackage{amsmath}
\usepackage{amssymb}
\usepackage{graphicx}
\usepackage{eufrak}
\usepackage{color}
\usepackage{slashed}
\usepackage{simplewick}
\usepackage{tikz}
\usepackage{tcolorbox}
\usepackage{ulem}
\graphicspath{{../figures/}}
%%figures
\def\lfig#1#2{\includegraphics[width=#1 in]{#2}}
\def\tfig#1#2{\includegraphics[height=#1 in]{#2}}
\def\addfig#1#2{\begin{center}\includegraphics[width=#1 in]{#2}\end{center}}
\def\question{\includegraphics[width=0.3in]{why.jpg}\,}
\def\answer{\includegraphics[width=0.3in]{baozou_haha.png}\,}
\def\wulian{\includegraphics[width=0.18in]{emoji_wulian.jpg}}
\def\bigwulian{\includegraphics[width=0.35in]{emoji_wulian.jpg}}
\def\bye{\includegraphics[width=0.18in]{emoji_bye.jpg}}
\def\bigbye{\includegraphics[width=0.35in]{emoji_bye.jpg}}
\def\huaixiao{\includegraphics[width=0.18in]{emoji_huaixiao.jpg}}
\def\bighuaixiao{\includegraphics[width=0.35in]{emoji_huaixiao.jpg}}
\def\jianxiao{\includegraphics[width=0.18in]{emoji_jianxiao.jpg}}
\def\bigjianxiao{\includegraphics[width=0.35in]{emoji_jianxiao.jpg}}
\def\haoqi{\includegraphics[width=0.18in]{emoji_haoqi.jpg}}
%% colors
\def\blacktext#1{{\color{black}#1}}
\def\bluetext#1{{\color{blue}#1}}
\def\redtext#1{{\color{red}#1}}
\def\darkbluetext#1{{\color[rgb]{0,0.2,0.6}#1}}
\def\skybluetext#1{{\color[rgb]{0.2,0.7,1.}#1}}
\def\cyantext#1{{\color[rgb]{0.,0.5,0.5}#1}}
\def\greentext#1{{\color[rgb]{0,0.7,0.1}#1}}
\def\darkgray{\color[rgb]{0.2,0.2,0.2}}
\def\lightgray{\color[rgb]{0.6,0.6,0.6}}
\def\gray{\color[rgb]{0.4,0.4,0.4}}
\def\blue{\color{blue}}
\def\red{\color{red}}
\def\orange{\color[rgb]{1.,0.8,0.}}
\def\green{\color{green}}
\def\darkgreen{\color[rgb]{0,0.4,0.1}}
\def\darkblue{\color[rgb]{0,0.2,0.6}}
\def\skyblue{\color[rgb]{0.2,0.7,1.}}
%%control
\def\diag{\mathrm{diag}\,}
\def\heaviside{\,\mathrm{h}\,}
\def\bral{\left(\begin{array}{l}}
\def\brar{\end{array}\right)}
\def\brall{\left(\begin{array}{ll}}
\def\brarr{\end{array}\right)}
\def\bralll{\left(\begin{array}{lll}}
\def\brarrr{\end{array}\right)}
\def\branchl{\left\{\begin{array}{l}}
\def\branchr{\end{array}\right.}
\def\branchll{\left\{\begin{array}{ll}}
\def\branchrr{\end{array}\right.}
\def\branchlll{\left\{\begin{array}{lll}}
\def\branchrrr{\end{array}\right.}
\def\sfgamma#1{\,\Gamma\left( #1 \right)\,}
\def\be{\begin{equation}}
\def\ee{\nonumber\end{equation}}
\def\bea{\begin{eqnarray}}
\def\eea{\nonumber\end{eqnarray}}
\def\bch{\begin{CJK}{UTF8}{gbsn}}
\def\ech{\end{CJK}}
\def\bitem{\begin{itemize}}
\def\eitem{\end{itemize}}
\def\bcenter{\begin{center}}
\def\ecenter{\end{center}}
\def\bex{\begin{minipage}{0.2\textwidth}\includegraphics[width=0.6in]{jugelizi.png}\end{minipage}\begin{minipage}{0.76\textwidth}}
\def\eex{\end{minipage}}
\def\chtitle#1{\frametitle{\bch#1\ech}}
\def\bmat#1{\left(\begin{array}{#1}}
\def\emat{\end{array}\right)}
\def\bcase#1{\left\{\begin{array}{#1}}
\def\ecase{\end{array}\right.}
\def\bmini#1{\begin{minipage}{#1\textwidth}}
\def\emini{\end{minipage}}
\def\tbox#1{\begin{tcolorbox}#1\end{tcolorbox}}
\def\pfrac#1#2#3{\left(\frac{\partial #1}{\partial #2}\right)_{#3}}
\def\res#1#2{\,\mathrm{res}\,#1\left(#2\right)\,}
\def\newt#1#2{\left(\begin{array}{c}#1\\ #2\end{array}\right)}
\def\reof#1{\,\mathrm{Re}{\left(#1\right)\,}}
\def\imof#1{\,\mathrm{Im}{\left(#1\right)\,}}
\def\Arg#1{\,\mathrm{Arg}\,#1\,}
%%symbols
\def\intfull{\int_{-\infty}^\infty}
\def\inthalf{\int_0^\infty}
\def\bropt{\,(\ \ \ )}
\def\sech{\mathrm{sech}\,}
\def\csch{\mathrm{csch}\,}
\def\asinh{\mathrm{asinh}\,}
\def\acosh{\mathrm{acosh}\,}
\def\sone{$\star$}
\def\stwo{$\star\star$}
\def\sthree{$\star\star\star$}
\def\sfour{$\star\star\star\star$}
\def\sfive{$\star\star\star\star\star$}
\def\rint{{\int_\leftrightarrow}}
\def\roint{{\oint_\leftrightarrow}}
\def\stdHf{{\textit{\r H}_f}}
\def\deltaH{{\Delta \textit{\r H}}}
\def\ii{{\dot{\imath}}}
\def\skipline{{\vskip0.1in}}
\def\skiplines{{\vskip0.2in}}
\def\lagr{{\mathcal{L}}}
\def\hamil{{\mathcal{H}}}
\def\vecT{{\mathbf{T}}}
\def\vecN{{\mathbf{N}}}
\def\vecB{{\mathbf{B}}}
\def\vecv{{\mathbf{v}}}
\def\vecr{{\mathbf{r}}}
\def\vecf{{\mathbf{f}}}
\def\vecg{{\mathbf{g}}}
\def\vecupsilon{{\mathbf{\upsilon}}}
\def\vecu{{\mathbf{u}}}
\def\vecj{{\mathbf{j}}}
\def\vecx{{\mathbf{x}}}
\def\vecy{{\mathbf{y}}}
\def\vecz{{\mathbf{z}}}
\def\veck{{\mathbf{k}}}
\def\vecp{{\mathbf{p}}}
\def\vecn{{\mathbf{n}}}
\def\vecA{{\mathbf{A}}}
\def\vecP{{\mathbf{P}}}
\def\vecsigma{{\mathbf{\sigma}}}
\def\hatJn{{\hat{J_\vecn}}}
\def\hatJx{{\hat{J_x}}}
\def\hatJy{{\hat{J_y}}}
\def\hatJz{{\hat{J_z}}}
\def\hatj#1{\hat{J_{#1}}}
\def\hatphi{{\hat{\phi}}}
\def\hatq{{\hat{q}}}
\def\hatpi{{\hat{\pi}}}
\def\vel{\upsilon}
\def\Dint{{\mathcal{D}}}
\def\adag{{\hat{a}^\dagger}}
\def\bdag{{\hat{b}^\dagger}}
\def\cdag{{\hat{c}^\dagger}}
\def\ddag{{\hat{d}^\dagger}}
\def\hata{{\hat{a}}}
\def\hatb{{\hat{b}}}
\def\hatc{{\hat{c}}}
\def\hatd{{\hat{d}}}
\def\hatD{{\,\hat{D}}}
\def\hatN{{\hat{N}}}
\def\hatH{{\hat{H}}}
\def\hatp{{\hat{p}}}
\def\Fup{{F^{\mu\nu}}}
\def\Fdown{{F_{\mu\nu}}}
\def\newl{\nonumber \\}
\def\vece{\mathrm{e}}
\def\calM{{\mathcal{M}}}
\def\calT{{\mathcal{T}}}
\def\calR{{\mathcal{R}}}
\def\barpsi{\bar{\psi}}
\def\baru{\bar{u}}
\def\barv{\bar{\upsilon}}
\def\qeq{\stackrel{?}{=}}
\def\ftf{\stackrel{\mathcal{FT}}{\Longrightarrow}}
\def\ftb{\stackrel{\mathcal{FT}}{\Longleftarrow}}
\def\ftfb{\stackrel{\mathcal{FT}}{\Longleftrightarrow}}
\def\ltf{\stackrel{\mathcal{LT}}{\Longrightarrow}}
\def\ltb{\stackrel{\mathcal{LT}}{\Longleftarrow}}
\def\ltfb{\stackrel{\mathcal{LT}}{\Longleftrightarrow}}
\def\torder#1{\mathcal{T}\left(#1\right)}
\def\rorder#1{\mathcal{R}\left(#1\right)}
\def\contr#1#2{\contraction{}{#1}{}{#2}#1#2}
\def\trof#1{\mathrm{Tr}\left(#1\right)}
\def\trace{\mathrm{Tr}}
\def\comm#1{\ \ \ \left(\mathrm{used}\ #1\right)}
\def\tcomm#1{\ \ \ (\text{#1})}
\def\slp{\slashed{p}}
\def\slk{\slashed{k}}
\def\calp{{\mathfrak{p}}}
\def\veccalp{\mathbf{\mathfrak{p}}}
\def\Tthree{T_{\tiny \textcircled{3}}}
\def\pthree{p_{\tiny \textcircled{3}}}
\def\dbar{{\,\mathchar'26\mkern-12mu d}}
\def\erf{\mathrm{erf}}
\def\const{\mathrm{const.}}
\def\pheat{\pfrac p{\ln T}V}
\def\vheat{\pfrac V{\ln T}p}
%%units
\def\fdeg{{^\circ \mathrm{F}}}
\def\cdeg{^\circ \mathrm{C}}
\def\atm{\,\mathrm{atm}}
\def\angstrom{\,\text{\AA}}
\def\SIL{\,\mathrm{L}}
\def\SIkm{\,\mathrm{km}}
\def\SIyr{\,\mathrm{yr}}
\def\SIGyr{\,\mathrm{Gyr}}
\def\SIV{\,\mathrm{V}}
\def\SImV{\,\mathrm{mV}}
\def\SIeV{\,\mathrm{eV}}
\def\SIkeV{\,\mathrm{keV}}
\def\SIMeV{\,\mathrm{MeV}}
\def\SIGeV{\,\mathrm{GeV}}
\def\SIcal{\,\mathrm{cal}}
\def\SIkcal{\,\mathrm{kcal}}
\def\SImol{\,\mathrm{mol}}
\def\SIN{\,\mathrm{N}}
\def\SIHz{\,\mathrm{Hz}}
\def\SIm{\,\mathrm{m}}
\def\SIcm{\,\mathrm{cm}}
\def\SIfm{\,\mathrm{fm}}
\def\SImm{\,\mathrm{mm}}
\def\SInm{\,\mathrm{nm}}
\def\SImum{\,\mathrm{\mu m}}
\def\SIJ{\,\mathrm{J}}
\def\SIW{\,\mathrm{W}}
\def\SIkJ{\,\mathrm{kJ}}
\def\SIs{\,\mathrm{s}}
\def\SIkg{\,\mathrm{kg}}
\def\SIg{\,\mathrm{g}}
\def\SIK{\,\mathrm{K}}
\def\SImmHg{\,\mathrm{mmHg}}
\def\SIPa{\,\mathrm{Pa}}
%page
\def\secpage#1#2{\begin{frame}\bcenter{\bf \Huge #1} \ecenter \skipline \tbox{\bcenter #2 \ecenter}\end{frame}}
\def\append#1#2{\secpage{附录#1}{#2}}
\def\thinka#1{\begin{frame}\frametitle{思考题}\bcenter\lfig{0.4}{think0.jpg}  \skipline #1 \ecenter\end{frame}}
\def\thinkb#1{\begin{frame}\frametitle{思考题}\bcenter\lfig{0.5}{think1.jpg}  \skipline #1 \ecenter\end{frame}}
\def\thinkc#1{\begin{frame}\frametitle{思考题}\bcenter\lfig{0.5}{think2.jpg}  \skipline #1 \ecenter\end{frame}}
\def\thinkd#1{\begin{frame}\frametitle{思考题}\bcenter\lfig{0.5}{think3.jpg}  \skipline #1 \ecenter\end{frame}}
\def\thinke#1{\begin{frame}\frametitle{思考题}\bcenter\lfig{0.5}{think4.jpg}  \skipline #1 \ecenter\end{frame}}
\def\thinkf#1{\begin{frame}\frametitle{思考题}\bcenter\lfig{0.5}{think5.jpg}  \skipline #1 \ecenter\end{frame}}
\def\schw{ds^2 = \left(1-\frac{2GM}{r}\right)dt^2 - \left(1-\frac{2GM}{r}\right)^{-1} dr^2 - r^2\left(d\theta^2 + \sin^2\theta d\phi^2\right)}

\input{titlepage.tex}
  \date{}
  \begin{document}
  \bch
  \tpage{17}{Schwarzschild Black Hole}


  \secpage{史瓦西黑洞的视界}{$$r=2GM$$}
  \begin{frame}
    对史瓦西度规，
    $$\schw $$
    在 $r<2GM$ 的区域内，其实 $r$ 坐标才是“时间”。如果 $r$ 减小代表时间增加，则称之为黑洞。如果 $r$  增大代表时间增加，则称之为白洞。

    \skipline

    （如果不谈量子效应）在黑洞 $r<2GM$ 区域里 $r$ 只能单调减少，任何物质无法从 $r<2GM$ 的区域跑出来。在白洞 $r<2GM$ 区域里 $r$ 只能单调增加，任何物质都无法进入 $r<2GM$ 的区域。

    \skipline

    因此，我们把史瓦西黑洞或者白洞的 {\blue $r=2GM$ }这个面称为\sout{结界}{\blue 视界}。  
  \end{frame}


  \begin{frame}
    一般恒星坍缩都是物质掉往中心的情况，所以是形成黑洞。目前天文观测上还没有任何迹象表明我们的可观测宇宙内存在白洞。

    \addfig{2}{bhwh.jpg}

    一种脑洞大开的想法是：如果白洞存在，它可能和黑洞在时空奇点相连，是旅行到“其他宇宙”的通道。
  \end{frame}


  \secpage{黑洞旅行}{有去无回……}
  
  \begin{frame}
    在之前讨论史瓦西度规里的粒子运动都是在 $r\gg GM$ 情况下讨论的。

    \addfig{1}{travel.jpg}

    为了研究 $r\sim GM$ 时的情况，我得去史瓦西黑洞旅行一趟。    
    
  \end{frame}


  \begin{frame}
    \frametitle{行李清单}
    \bmini{0.5}
    \bitem
  \item[]{
    
    \tfig{0.8}{clock.jpg}
  }
  \item[]{    
    \tfig{0.8}{laserpointer.jpg}
  }
  \item[]{    
    \tfig{1}{RiemannTensorMachine.jpg}
  }
    \eitem
    \emini
    \bmini{0.45}
    \bitem
  \item{钟：计时。    \vspace{0.6in}
  }
  \item{激光笔：向空间站传输信息。\vspace{0.6in}}
  \item{黎曼曲率测量仪：测量时空曲率，确定自己位置。}
    \eitem
    \emini
  \end{frame}

  \begin{frame}
    假设在$r=r_0$ ($r_0\gg GM$) 的地方有相对史瓦西坐标保持静止($r,\theta,\phi$不变）的空间站。这就是我的出发点了。

      \addfig{1}{chufa.jpg}

    对径向下落者，不能使用Binet方程，要重新进行推导:
  \end{frame}


  \secpage{钟的计时}{下落时间有限且有解析解}
  
  \begin{frame}

    假设我从空间站开始由静止下落，则初始时刻，由归一化条件得到我的四维速度协变分量：
    $$ u_t = \sqrt{1-\frac{2GM}{r_0}} $$
    由于 $u_t$ 守恒，在任意时刻有：
    $$ \frac{dt}{ds} = \frac{\sqrt{1-\frac{2GM}{r_0}}}{1-\frac{2GM}{r}}$$
    然后利用 $ds^2$ 的表达式，有
    $$ \left(1-\frac{2GM}{r}\right)\left(\frac{dt}{ds}\right)^2 - \left(1-\frac{2GM}{r}\right)^{-1}\left(\frac{dr}{ds}\right)^2 = 1$$
    由此联立解出
    $$\frac{dr}{ds} = -\sqrt{\frac{2GM}{r}-\frac{2GM}{r_0}}$$
  \end{frame}


  \begin{frame}
    $$\frac{dr}{ds} = -\sqrt{\frac{2GM}{r}-\frac{2GM}{r_0}}$$
    这方程和牛顿力学自由下落是一样的！ 我们熟知牛顿力学里对应问题的解为滚轮线。于是照搬得到：

    \bea
    r &=& \frac{r_0}{2}(1-\cos\theta) \newl
    s &=& \sqrt{\frac{r_0^3}{8GM}}\left(\theta-\sin\theta\right)
    \eea
    这里的 $\theta\in[\pi, 2\pi]$。初始时刻对应滚轮线顶点 $\theta = \pi$，掉到奇点的时刻对应 $\theta=2\pi$。

  \end{frame}
  

  \begin{frame}
    可见，在掉到奇点前我的钟走过的时间是有限的：
    $$ s(2\pi)-s(\pi)  = \frac{\pi}{2}\sqrt{\frac{r_0^3}{2GM}}$$
  \end{frame}


  \secpage{激光笔通信}{多普勒红移+引力红移}
  
  \begin{frame}
    在掉进视界之后我当然就无法用激光笔和空间站联络了。在此之前，我的激光笔发出的光子，在我的局域 Minkowski 标架下观察的频率为固定的 $\nu_0$（因为我在自由下落；对非自由下落运动者而言这件事其实还蛮伤脑筋的）。
    即把光子的四维动量投影到我的局域标架时间轴（即我的四维速度 $u_\mu$）上，满足
    $$p_t \frac{\sqrt{1-\frac{2GM}{r_0}}}{1-\frac{2GM}{r}} + p_r \sqrt{\frac{2GM}{r}-\frac{2GM}{r_0}}  = h\nu_0$$
    又光子本身走零测地线，所以
    $$ (p_t)^2\left(1-\frac{2GM}{r}\right)^{-1} - (p_r)^2\left(1-\frac{2GM}{r}\right) = 0 $$
    联立上面两式可以解出 
  \end{frame}
  
  \begin{frame}
    $$p_t = h\nu_0\left(\sqrt{1-\frac{2GM}{r_0}} - \sqrt{\frac{2GM}{r}-\frac{2GM}{r_0}}\right)$$
    由于光子的 $p_t$ 守恒，空间站那里接收到的光子能量为（即把光子四维动量投影到空间站的四维速度上）：
    $$ h\nu' = p_t \frac{1}{\sqrt{1-\frac{2GM}{r_0}}} = h\nu_0\left(1 - \sqrt{\frac{\frac{2GM}{r}-\frac{2GM}{r_0}}{1-\frac{2GM}{r_0}} }\right) $$
    即当 $r$ 无限接近视界 $2GM$ 时，我发出的光子全都被无限红移了。这些光子到达空间站需要经过的时间也越来越长，趋于无穷。

  \end{frame}

  \thinkc{如果一个相对于史瓦西坐标系静止的人在 $r<r_0$ 处给空间站发射光子，引力红移效应是更强还是更弱？}

  \secpage{我眼中的天空}{逐渐消失的圆盘}

  \begin{frame}
    旅行中的我突然想用激光笔向尽可能多的方向发出信号（毕竟马上快失联了，总想留下点什么）

    只要光不沿径向走，我们就可以用之前推导的光子版Binet方程：
    $$\frac{d^2u}{d\phi^2}+u = 3u^2$$
    这里的 $u=\frac{GM}{r}$。

    设初始条件为 $u=u_0<\frac{1}{2}$，我需要计算初始 $\frac{du}{d\phi}$ 可能的取值范围(这里不妨假设我选择了 $\phi$ 单调增加的方向发射光子)，使得发出的光能跑到无穷远处（即 $u\rightarrow 0$)。
  \end{frame}


  \begin{frame}
    \frametitle{奇怪弹簧}
    \bmini{0.45}
    假设我有个质量为$1$，弹簧势能函数为 $V(u) = \frac{1}{2}u^2-u^3$ 的奇怪振子，$u$ 是振子的位移（弹簧的伸长量）。

    那么振子的运动方程为：
    $$ \ddot u = -\frac{dV}{du} = -u + 3u^2 $$    
    \emini
    \bmini{0.5}
    \ \lfig{2.2}{Vofu.png}
    \emini
    

    于是问题转化为，已知初始位移 $u_0$ 的情况下，初始速度$\dot u$ 为多少才能使振子回到 $u=0$ 的位置？显然答案是：

    \bitem
  \item{当 $u_0<\frac{1}{3}$，要求初始 $\dot u < \sqrt{2\left[V\left(\frac{1}{3}\right)-V\left(u_0\right)\right]}$}
  \item{当 $u_0\ge \frac{1}{3}$，要求初始 $\dot u < - \sqrt{2\left[V\left(\frac{1}{3}\right)-V\left(u_0\right)\right]}$}
    \eitem
  \end{frame}

  \begin{frame}
    对应回黑洞的问题，
    \bitem
  \item{当$u_0<\frac{1}{3}$ 即 $r>3GM$ 时，要求初始 $\frac{du}{d\phi} < \sqrt{\frac{1}{27}-u_0^2(1-2u_0)}$}
  \item{当 $\frac{1}{3}<u_0\le \frac{1}{2}$即 $2GM<r\le 3GM$时，要求初始 $\frac{du}{d\phi}<-\sqrt{\frac{1}{27}-u_0^2(1-2u_0)}$}
  \item{当 $r<=2GM$ 时，对应让 $u$ 减小（$r$ 增加）的解是白洞，和黑洞的假设不符，故要舍去这种可能性。}
    \eitem
  \end{frame}

  \begin{frame}
    继续把初始的 $\frac{du}{d\phi}$ 转化为我看到的出射角（这需要把光子的四维速度往我的局域标架上投影）这件事情没有多大的乐趣。没乐趣的事情就留给你做了……（？？？）

    \skipline
    
    可以看到的是，随着我靠近视界面，光线的出射方向范围越来越窄。那么，对于进来的光线也是如此。我能看到的天空逐渐缩小，最后在我到达黑洞视界处时完全消失（这一点你光看 $du/d\phi$ 并不能看出来，需要做投影计算）。然后，我进入了黑洞视界……

    \addfig{1.5}{fengxiaoxiao.jpg}
  \end{frame}


  \ech
\end{document}




  
